\documentclass[10pt,a4paper]{amsart}
\usepackage[margin=19mm]{geometry}
\usepackage{amsmath,amssymb,mathtools}
\usepackage[scr=rsfs,cal=euler]{mathalfa}
\usepackage{xfrac}
\usepackage[unicode,hidelinks,bookmarks=false]{hyperref}
\newcommand{\ie}{i.e.\ }
\newcommand{\cf}{cf.\ }
\newcommand{\resp}{respectively\ }
\newcommand{\Subsection}{\subsection}
\providecommand{\nobreakdash}{\nobreak\hbox{-}}
\setlength{\emergencystretch}{2em}
\multlinegap0pt
\usepackage{amssymb}
\usepackage{tikz}
\newtheorem{prop}{Proposition}
\def\Max{{\rm Max}}
\def\idem{{\rm idem}}
\def\nil{{\rm nil}}
\def\u{\mathrm{u}}
\def\vnr{{\rm vnr}}
\def\z{{zero}}
\title{On the $S$-version of some special elements in commutative rings}
\author{D. Bennis, A. Bouziri, S. D. Kumar,, T. Singh}
\date{Source: arXiv:2603.18890v1 (CC BY 4.0)}
\begin{document}
\maketitle

\noindent\emph{Source and licence.} On the $S$-version of some special elements in commutative rings. Version arXiv:2603.18890v1 (CC BY 4.0), distributed by arXiv under the Creative Commons Attribution 4.0 International licence, http://creativecommons.org/licenses/by/4.0/, as recorded in the arXiv licence metadata for this exact version and shown on the abstract page https://arxiv.org/abs/2603.18890v1. Full licence notice: \texttt{licenses/arxiv-2603.18890-CC-BY-4.0.txt}.\par\smallskip

\noindent\textit{Editorial scope.} Counted unit: the article's prop prop (source0.tex lines 211--220). The article's complete original proof of it is delivered with it (source0.tex lines 222--276, one \texttt{proof} environment of 55 lines). No mathematics is written, repaired or re-derived: the statement and the proof are the article's own lines, in the article's own notation, and no retained line is substituted or re-flowed.  No further source-local context is needed: every cross-reference of every retained line resolves inside the two delivered spans.  Nothing was omitted from any retained span, so the package carries no omission record.  No retained line contains a citation, so the wrapper carries no bibliography and no bibliography entry.  No retained line contains a figure environment, an image-inclusion command or drawing code, so no image is delivered and the package declares no asset dependency.  Editorial formatting only, in addition: an AMS \texttt{amsart} preamble that restates the article's own declarations and macros used by the retained lines, namely the environments \texttt{prop} on separate counters, together with the standard packages those lines need.  The retained environments print in this document as Proposition 1 in order of appearance, whereas the article prints the same items with the numbering of its own full document.  The complete native source of the article as distributed in the arXiv e-print for this exact version is bundled unchanged as \texttt{sources/196734/source0.tex}.
\begin{prop}
The following assertions hold:
\begin{enumerate}
    \item \( \{0\} = \bigcap_{M \in \Max(R)} M\text{-}\z(R) \);
    \item \( \u(R) = \bigcap_{M \in \Max(R)} M\text{-}\u(R) \);
    \item \( \nil(R) = \bigcap_{M \in \Max(R)} M\text{-}\nil(R) \);
    \item \( \idem(R) \subseteq \bigcap_{M \in \Max(R)} M\text{-}\idem(R) \subseteq \vnr(R) \);
    \item \( \vnr(R) = \bigcap_{M \in \Max(R)} M\text{-}\vnr(R) \).
\end{enumerate}
\end{prop}

\begin{proof}
(1)  Let \( a \in \bigcap_{M \in \Max(R)} M\text{-}\z(R) \). For each maximal ideal \( M \), there exists an element \( t_M \notin M \) such that \( t_M a = 0 \). Since the ideal generated by the family \( \{ t_M : M \in \Max(R) \} \) is equal to \( R \), there exist elements \( r_1, \dots, r_n \in R \) and maximal ideals \( M_1, \dots, M_n \in \Max(R) \) such that
\[
1 = r_1 t_{M_1} + \cdots + r_n t_{M_n}.
\]
It follows that
\[
a = r_1 t_{M_1} a + \cdots + r_n t_{M_n} a = 0.
\]

(2) Let \( a \in \bigcap_{M \in \Max(R)} M\text{-}\u(R) \). For each maximal ideal \( M \), there exist \( t_M \notin M \) and \( x_M \in R \) such that \( a x_M = t_M \). As in (1), the family \(\{ t_M : M \in \Max(R) \} \) generate the unit ideal, so
\[
1 = r_1 t_{M_1} + \cdots + r_n t_{M_n}
\]
for some \( r_i \in R \) and \( M_i \in \Max(R) \). Hence,
\[
1 = r_1 a x_{M_1} + \cdots + r_n a x_{M_n} = a(r_1 x_{M_1} + \cdots + r_n x_{M_n}),
\]
which shows that \( a \in \u(R) \).

(3) Let \( a \in \bigcap_{M \in \Max(R)} M\text{-}\nil(R) \). For each maximal ideal \( M \), there exists \( t_M \notin M \) and \( p_M \in \mathbb{N} \) such that \( t_M a^{p_{M}} = 0 \). Since  the family \(\{ t_M : M \in \Max(R) \} \) generate the unit ideal, there exist \( r_1, \dots, r_n \in R \) and \( M_1, \dots, M_n \in \Max(R) \) such that
\[
1 = r_1 t_{M_1} + \cdots + r_n t_{M_n}.
\]
Set \( N = p_{M_1} + \cdots + p_{M_n} \). Then,
\[
a^N = a^N (r_1 t_{M_1} + \cdots + r_n t_{M_n}) = \sum r_i a^N t_{M_i} = 0,
\]
since each \( t_{M_i} a^{p_{M_i}} = 0 \), and \( a^N = a^{q} a^{p_{M_i}} \) for some \( q \ge 0 \). Hence \( a \in \nil(R) \).

(4) The first inclusion is clear. Let \( a \in \bigcap_{M \in \Max(R)} M\text{-}\idem(R) \). Then for each \( M \), there exists \( t_m \notin M \) such that \( t_M a = a^2 \). As above, we write
\[
1 = r_1 t_{M_1} + \cdots + r_n t_{M_n}.
\]
Then,
\[
a = r_1 t_{M_1} a + \cdots + r_n t_{M_n} a = r_1 a^2 + \cdots + r_n a^2 = (r_1 + \cdots + r_n) a^2,
\]
so \( a \in \vnr(R) \). If \( a \) is regular, then
\[
1 = (r_1 + \cdots + r_n) a.
\]

(5) This follows by using a similar argument as in the previous statements.

% Let \( a \in \bigcap_{m \in \Max(R)} m\text{-}\vnr(R) \). For each maximal ideal \( m \), there exist \( t_m \notin m \) and \( x_m \in R \) such that \( t_m a = x_m a^2 \). Again,
%\[
%1 = r_1 t_{m_1} + \cdots + r_n t_{m_n}.
%\]
%Then,
%\[
%a = r_1 t_{m_1} a + \cdots + r_n t_{m_n} a = r_1 x_{m_1} a^2 + \cdots + r_n x_{m_n} a^2 = (r_1 x_{m_1} + \cdots + r_n x_{m_n}) a^2,
%\]
%which shows that \( a \in \vnr(R) \).
\end{proof}

\end{document}
